
\section{Model formulation details}
\label{app:model_details}

\subsection{Heat-pump performance pre-computation}
\label{app:cop}

The waste-heat coefficient of performance is pre-computed from an analytical Lorenz-cycle
model,
\begin{equation}
  \mathrm{COP}_{h,t}^{\mathrm{WHR}} = \eta_{\mathrm{Lorenz}}
  \cdot \frac{T_{\mathrm{sink,out}}}{T_{\mathrm{sink,out}} -
  T_{\mathrm{src},h}(t) + \Delta T_{\mathrm{pp}}},
\end{equation}
with all temperatures in kelvin, a Lorenz efficiency of $0.6$ and a pinch-point difference of $5$\,K, giving a combined
approximation error of about $3.6$\,\%. Pre-computing rather than co-optimising the
coefficient of performance keeps the formulation free of bilinear terms; it also decouples
heat-pump operation from the network state, which is one of the assumptions discussed in
Section~\ref{subsec:whynull} as limiting how far the extended physics can move a dispatch
decision.

\subsection{Linear component models}
\label{app:component_models}

{%
\captionof{table}{Linear component models, applied identically at every fidelity level.}
\label{tab:simple_components}
\centering
\footnotesize
\begin{tabular*}{\columnwidth}{@{\extracolsep{\fill}} l l l @{}}
\toprule
Component & Equation & Bound \\
\midrule
Gas boiler       & $Q_{g,t}^{\mathrm{th}} = \eta_g F_{g,t}$ & $\leq \bar{Q}_g$ \\
Biomass boiler   & $Q_{g,t}^{\mathrm{th}} = \eta_g F_{g,t}$ & $\leq \bar{Q}_g$ \\
Electrode boiler & $Q_{r,t} = \eta_r P_{r,t}^{\mathrm{el}}$ & $\leq \bar{Q}_r$ \\
CHP              & $Q^{\mathrm{th}} = \eta^{\mathrm{th}} F$ & $\sigma_g$ constant \\
\bottomrule
\end{tabular*}
}% end centering group

\bigskip

Because these models are identical at every level, they cannot contribute to any
level-to-level difference. This is what allows the contrasts of
Table~\ref{tab:contrasts} to be attributed to the network representation alone.

\subsection{Piecewise-linearisation error bound}
\label{app:pwl}

For the quadratic pressure relation $\Delta p = R Q^2$ approximated over $K$ equal
segments, the maximum approximation error is
\begin{equation}
  \varepsilon_{\Delta p}^{\max} = \frac{R \bar{Q}^2}{4K^2},
\end{equation}
which at $K = 3$ is $2.8$\,\% of the peak pressure drop. The resulting cost deviation is far
smaller than this bound, because the optimiser compensates by adjusting the schedule rather
than by violating feasibility: the approximation error manifests as a minor scheduling
shift, not as a proportional cost deviation. The cost deviation itself is not estimated from
this bound but measured directly, by re-solving the continuous problem under native
non-linearised physics with the mixed-integer schedule fixed
(Section~\ref{subsec:linearisation}).

\medskip
{%
\captionof{table}{Piecewise-linear approximation error and problem size against the number
of segments $K$. The cost effect is not inferred from this table; it is measured by the
solved nonlinear reference in Section~\ref{subsec:linearisation}.}
\label{tab:pwl_sensitivity}
\centering
\footnotesize
\begin{threeparttable}
\begin{tabular*}{\columnwidth}{@{\extracolsep{\fill}} c c r @{}}
\toprule
$K$ & $\varepsilon^{\max}$ & Additional variables\tnote{a} \\
\midrule
3 & 2.8\,\% & 753k \\
5 & 1.0\,\% & 998k\tnote{b} \\
8 & 0.4\,\% & 1366k\tnote{b} \\
\bottomrule
\end{tabular*}
\begin{tablenotes}\scriptsize
  \item[a] Relative to the node-resolved baseline, binary and continuous combined.
  \item[b] Extrapolated at approximately 122k variables per additional segment; not
    simulated.
\end{tablenotes}
\end{threeparttable}
}% end centering group

\subsection{Taylor linearisation of the temperature product}
\label{app:taylor_derivation}

The bilinear product of upstream temperature and decay factor is expanded about the
nominal operating point $(T_{n_1}^0, \phi_p^0)$:
\begin{equation}
  T_{n_1} \phi_p \approx T_{n_1}^0 \phi_p^0
  + \phi_p^0 (T_{n_1} - T_{n_1}^0)
  + T_{n_1}^0 (\phi_p - \phi_p^0).
\end{equation}
The neglected second-order cross-term is bounded by
$|\varepsilon| \leq |\Delta T_{\max}| |\Delta\phi_{\max}| < 2.5$\,K over the operating
envelope. The exponential decay factor is approximated piecewise-linearly only in the
PWL-linearised temperature-propagating level ($\Ltwo$); the nonlinear reference treats it
natively (Section~\ref{subsec:evaluator}). The measured linearisation gap of
$-0.15$/$-0.33$\,\% (Section~\ref{subsec:linearisation}) therefore already contains this
exponential-approximation component (native versus piecewise-linear) rather than cancelling
it, so the reported figure is the combined PWL-versus-native gap, not an isolated one.
